\documentclass[10pt,a4paper]{amsart}
\usepackage[margin=19mm]{geometry}
\usepackage{amsmath,amssymb,mathtools}
\usepackage[scr=rsfs,cal=euler]{mathalfa}
\usepackage{xfrac}
\usepackage[unicode,hidelinks,bookmarks=false]{hyperref}
\newcommand{\ie}{i.e.\ }
\newcommand{\cf}{cf.\ }
\newcommand{\resp}{respectively\ }
\newcommand{\Subsection}{\subsection}
\providecommand{\nobreakdash}{\nobreak\hbox{-}}
\setlength{\emergencystretch}{2em}
\multlinegap0pt
\usepackage{booktabs}
\usepackage{mathtools}
\usepackage[shortlabels]{enumitem}
\usepackage{xcolor}
\usepackage{algorithm}\usepackage{algpseudocode}
\usepackage{amssymb}
\newtheorem{proposition}{Proposition}
\providecommand{\R}{\mathbb{R}}
\providecommand{\diag}{\mathrm{diag}}
\DeclareMathOperator{\rank}{rank}
\title{Beyond Low Rank: Fast Low-Rank + Diagonal Decomposition with a Spectral Approach}
\author{Kingsley Yeon, Mihai Anitescu}
\date{Source: arXiv:2512.17120v1 (CC BY 4.0)}
\begin{document}
\maketitle

\noindent\emph{Source and licence.} Beyond Low Rank: Fast Low-Rank + Diagonal Decomposition with a Spectral Approach. Version arXiv:2512.17120v1 (CC BY 4.0), distributed by arXiv under the Creative Commons Attribution 4.0 International licence, http://creativecommons.org/licenses/by/4.0/, as recorded in the arXiv licence metadata for this exact version and shown on the abstract page https://arxiv.org/abs/2512.17120v1. Full licence notice: \texttt{licenses/arxiv-2512.17120-CC-BY-4.0.txt}.\par\smallskip

\noindent\textit{Editorial scope.} Counted unit: the article's proposition prop:naivedecomposition (source0.tex lines 146--178). The article's complete original proof of it is delivered with it (source0.tex lines 180--218, one \texttt{proof} environment of 39 lines). No mathematics is written, repaired or re-derived: the statement and the proof are the article's own lines, in the article's own notation, and no retained line is substituted or re-flowed.  No further source-local context is needed: every cross-reference of every retained line resolves inside the two delivered spans.  Nothing was omitted from any retained span, so the package carries no omission record.  No retained line contains a citation, so the wrapper carries no bibliography and no bibliography entry.  No retained line contains a figure environment, an image-inclusion command or drawing code, so no image is delivered and the package declares no asset dependency.  Editorial formatting only, in addition: an AMS \texttt{amsart} preamble that restates the article's own declarations and macros used by the retained lines, namely the environments \texttt{proposition} on separate counters, together with the standard packages those lines need.  The retained environments print in this document as Proposition 1 in order of appearance, whereas the article prints the same items with the numbering of its own full document.  The complete native source of the article as distributed in the arXiv e-print for this exact version is bundled unchanged as \texttt{sources/191293/source0.tex}.
\begin{proposition}[Existence and Error Reduction in a Na\"ive Diagonal + Low-Rank Approximation] \label{prop:naivedecomposition}
Let \(\Sigma\in\R^{n\times n}\) be symmetric positive semidefinite with eigen-decomposition
\[
  \Sigma = V\Lambda V^T,\quad
  \Lambda=\operatorname{diag}(\lambda_1,\dots,\lambda_n),\quad
  \lambda_1\ge\cdots\ge\lambda_n\ge0.
\]
Fix \(0\le k\le n\), and write \(V=[V_k\;\;V_{>k}]\), \(\Lambda=\mathrm{diag}(\Lambda_k,\Lambda_{>k})\).  Define
\[
  U_k = V_k\,\Lambda_k^{1/2}, 
  \qquad
  D_k = \diag\!\bigl(\Sigma_{ii}-(U_kU_k^T)_{ii}\bigr),
\]
and set the residual
\[
  S_k = \Sigma - U_kU_k^T
  = V_{>k}\,\Lambda_{>k}\,V_{>k}^T \succeq 0,
\]
and the “corrected’’ residual
\[
  R_k = \Sigma - (D_k + U_kU_k^T)
  = S_k \;-\;D_k.
\]
Then:
\begin{enumerate}
  \item \(\Sigma = D_k + U_kU_k^T + R_k\), with \(D_k\succeq0\).
  \item \(\|S_k\|_2 = \lambda_{k+1}\) and
  \(\|R_k\|_2 \le \lambda_{k+1}\), with strict inequality
  \(\|R_k\|_2<\lambda_{k+1}\) whenever \(\lambda_{k+1}>0\).  
  \item If \(\lambda_{k+1}=0\) (i.e.\ \(\rank\Sigma\le k\)), then \(D_k=0\) and \(R_k=0\), so \(\Sigma=U_kU_k^T\).
\end{enumerate}
In particular, whenever \(\Sigma\) is full-rank beyond \(k\), including \(D_k\) strictly reduces the operator-norm error compared to the pure rank-\(k\) approximation.
\end{proposition}

\begin{proof}
By construction \(U_kU_k^T=V_k\Lambda_kV_k^T\), so 
\[
S_k=\Sigma-U_kU_k^T=V_{>k}\Lambda_{>k}V_{>k}^T,\quad
\|S_k\|_2=\lambda_{k+1}.
\]
Since \((U_kU_k^T)_{ii}\le\Sigma_{ii}\), we have \(D_k\succeq0\).  
Moreover, recall that \(R_k = S_k - \diag(S_k)\).  
To show \(\|R_k\|_2 \le \|S_k\|_2\), note that for any unit vector \(x\),
\[
|x^\top R_k x| = |x^\top S_k x - x^\top \diag(S_k) x|
  \le \max\{x^\top S_k x,\; x^\top \diag(S_k) x\}.
\]
Taking the supremum over all \(\|x\|_2=1\) gives
\[
\|R_k\|_2 = \sup_{\|x\|=1} |x^\top R_k x|
   \le \max\{\|S_k\|_2,\|\diag(S_k)\|_2\}.
\]
Since \(S_k \succeq 0\), the diagonal matrix \(D = \diag(S_k)\) is also positive semidefinite.  
For a diagonal matrix, its spectral norm equals its largest diagonal entry in absolute value, and since \(S_k\) is positive semidefinite, these entries are nonnegative. Hence
\[
\|D\|_2 = \max_i D_{ii} = \max_i (S_k)_{ii}.
\]
Each diagonal entry satisfies \((S_k)_{ii} = e_i^\top S_k e_i\), where \(e_i\) is the \(i\)-th standard basis vector.  
Because for any unit vector \(x\), \(x^\top S_k x \le \lambda_{\max}(S_k)\), we have
\[
(S_k)_{ii} = e_i^\top S_k e_i \le \lambda_{\max}(S_k) \quad \forall i.
\]
Taking the maximum over \(i\) yields
\[
\|D\|_2 = \max_i (S_k)_{ii} = \max_i e_i^\top S_k e_i \le \lambda_{\max}(S_k) = \|S_k\|_2.
\]
we conclude that
\[
\|R_k\|_2 \le \|S_k\|_2.
\]

Moreover, if \(\lambda_{k+1}>0\), then \(\diag(S_k)\neq0\), so subtracting it strictly lowers the Rayleigh quotient of the top eigenvector of \(S_k\) and hence \(\|R_k\|_2<\lambda_{k+1}\).  Finally, if \(\lambda_{k+1}=0\) then \(S_k=0\), forcing \(D_k=0\) and \(R_k=0\).
\end{proof}

\end{document}
